\documentclass[12pt]{article}
\usepackage[utf8]{inputenc}
\usepackage[T1]{fontenc}
\usepackage{hyperref}
\title{Kaimbulawa Language Resources for K-12 Curriculum Development}
\author{Compiled by Claude Sonnet 5.5 (claude-sonnet-5-5) for Andr\'as Kornai}
\date{}
\begin{document}
\maketitle

\section{Introduction}
Kaimbulawa (ISO 639-3: zka) is an Austronesian language of Siompu Island, off Buton Island in Southeast Sulawesi, Indonesia. It belongs to the Muna--Buton group. Wikipedia cites about 2,200 speakers (2005) and Global Recordings gives 2,800. The people live in four villages: Kaimbulawa, Waidawula, Lontoy and Karae. Only general reference pages turned up in the searches.

\section{Primary Educational Resources}
No primers, textbooks or school programmes for Kaimbulawa turned up in the searches.

\section{Linguistic and Reference Materials}
\begin{itemize}
\item Wikipedia, Kaimbulawa language (a stub article): \url{https://en.wikipedia.org/wiki/Kaimbulawa_language}
\end{itemize}

\section{Religious and Community Resources}
\begin{itemize}
\item Joshua Project, Kaimbulawa in Indonesia: \url{https://joshuaproject.net/people_groups/12425}
\item Asia Harvest people-group profile (PDF): \url{https://people-groups.asiaharvest.org/Indonesia/Kaimbulawa.pdf}
\end{itemize}

\section{Digital and Audio Resources}
\begin{itemize}
\item Global Recordings Network, Kaimbulawa language page (the listing of recordings was not checked): \url{https://globalrecordings.net/en/language/zka}
\end{itemize}

\section{Curriculum Development Context}
Little published material exists for Kaimbulawa. A curriculum would need new documentation, an orthography and teaching materials, ideally developed with the community. Indonesian is the language of schooling in Indonesia.
\end{document}